% report.tex -- compile with: latexmk -lualatex report.tex
\documentclass[11pt,a4paper]{article}

\usepackage[margin=2.2cm]{geometry}
\usepackage{amsmath,amssymb,amsthm}
\usepackage{fontspec}
\usepackage{unicode-math}
\setmainfont{STIX Two Text}
\setmathfont{STIX Two Math}
\setmonofont{DejaVu Sans Mono}[Scale=0.80]
\usepackage{booktabs,longtable,array}
\usepackage{fancyvrb}
\usepackage{luacode}
\usepackage{enumitem}
\usepackage[hidelinks]{hyperref}
\setlist{itemsep=1pt,topsep=3pt}
\setlength{\emergencystretch}{3em}

% ---------------------------------------------------------------- Lean names
% \lean{name}: typeset a Lean identifier verbatim in monospace, with break
% points after '_' and '.', so long names do not overflow table cells.
\begin{luacode*}
-- break after '_' '.' ':' freely; at a lower->upper camelCase boundary with a small penalty
function leanname(s)
  local buf = ""
  local prev = ""
  for _, c in utf8.codes(s) do
    local ch = utf8.char(c)
    if prev:match("^%l$") and ch:match("^%u$") then
      tex.sprint(-2, buf); buf = ""
      tex.sprint("\\penalty100\\hskip0pt\\relax{}")
    end
    buf = buf .. ch
    if ch == "_" or ch == "." or ch == ":" or ch == "/" then
      tex.sprint(-2, buf); buf = ""
      tex.sprint("\\allowbreak{}")
    end
    prev = ch
  end
  tex.sprint(-2, buf)
end
\end{luacode*}
\newcommand{\lean}[1]{\texttt{\directlua{leanname("\luaescapestring{\detokenize{#1}}")}}}
\newcommand{\lbl}[1]{\lean{#1}}   % paper labels (b-prefixed for Barbot)
\newcommand{\dec}[1]{DEC-#1}

% ---------------------------------------------------------------- maths
\newcommand{\R}{\mathbb{R}}
\newcommand{\Z}{\mathbb{Z}}
\newcommand{\N}{\mathbb{N}}
\newcommand{\C}{\mathbb{C}}
\newcommand{\Mink}{\mathbf{M}}
\newcommand{\Dev}{\mathcal{D}}
\newcommand{\Sc}{\widetilde{S}}
\newcommand{\Isom}{\operatorname{Isom}}
\newcommand{\SOo}{\mathrm{SO}_0}
\newcommand{\sing}{\operatorname{sing}}
\newcommand{\reg}{\operatorname{reg}}
\newcommand{\inter}{\operatorname{int}}

\theoremstyle{plain}
\newtheorem{theorem}{Theorem}[section]
\newtheorem{proposition}[theorem]{Proposition}
\theoremstyle{definition}
\newtheorem{definition}[theorem]{Definition}
\newtheorem{reading}[theorem]{Reading}
\theoremstyle{remark}
\newtheorem*{remark}{Remark}

